\title{LongRoPE: Extending LLM Context Window Beyond 2 Million Tokens}

\begin{abstract}
Expanding the context window in large language models (LLMs) is highly sought after. Nevertheless, limitations persist due to the substantial resources needed for fine-tuning, the limited availability of lengthy textual data, and problematic token positions that cause catastrophic values. As a result, modern LLMs typically support context windows up to about 128k tokens.

This work presents LongRoPE, a method that, for the first time, pushes the context window of pre-trained LLMs to an impressive \textbf{2048k} tokens. Remarkably, this extension requires as few as 1k fine-tuning steps, limited to a maximum context length of 256k during training, all while preserving performance at the original short context range. The approach hinges on three principal innovations: (i) we uncover and leverage two distinct types of non-uniformity inherent in positional interpolation through an efficient search process, yielding improved initialization for fine-tuning and enabling an 8$\times$ context length increase without further fine-tuning; (ii) we develop a stepwise extension approach, starting with the fine-tuning of an LLM on a 256k sequence length, followed by a subsequent positional interpolation applied to this already-extended model, thereby attaining a 2048k window; (iii) we further realign LongRoPE at the 8k length to restore short context window performance. Comprehensive experiments on LLaMA2 and Mistral spanning multiple tasks affirm the robustness of our method. Models upgraded through LongRoPE preserve the original architecture, requiring only minimal alterations to the positional embedding, and remain compatible with the vast majority of existing optimizations. Code will be made accessible at \texttt{https://github.com/microsoft/LongRoPE}

\section{Introduction}

Although Large Language Models (LLMs) have achieved impressive results across a range of tasks~\cite{gpt4,llama2}, they are constrained by a restricted context window size—LLaMA2, for instance, is limited to 4096 tokens~\cite{llama2}. Once inputs exceed this boundary, model effectiveness diminishes, largely because the architecture has not been trained to handle tokens in these additional positions. Such constraints hinder performance in key applications, including in-context learning with many examples~\cite{Huang2023FewerIM} as well as the development of LLM-based agents~\cite{park2023generative,madaan2023selfrefine}.

Recent studies demonstrate that by fine-tuning pre-trained LLMs with longer input sequences, the context window can be increased to approximately 128k tokens~\cite{longlora,pi,yarn,zhang2024soaring,liu2023scaling}. However, pushing the context window beyond this range faces three significant challenges. Firstly, introducing untrained positional indices generates numerous problematic values, resulting in distributional shifts that hinder the convergence of fine-tuning~\cite{pi}. This difficulty is exacerbated in cases where extending the context window from 4k to $>$1000k adds over 90\% new positions. Secondly, effective fine-tuning typically depends on having access to training samples that match the desired sequence lengths, yet available datasets rarely provide sufficiently long texts, particularly those surpassing 1000k tokens. Training with these extremely long sequences is also highly resource-intensive, demanding excessive computational time and substantial GPU allocation. Thirdly, when context windows are enlarged to extreme lengths, attention mechanisms become diluted across vast token positions, causing the model's effectiveness on shorter contexts—its original regime—to deteriorate~\cite{pi}.

A commonly used strategy to address the initial challenge involves interpolating the RoPE positional embedding~\cite{rope,pi}, which rescales new positional indices to align with the range seen during pretraining, as illustrated in Fig.\ref{fig:overview}. Position Interpolation (PI) ~\cite{pi} performs a linear interpolation of RoPE’s rotary angles, guided by the extension ratio. The NTK method~\cite{ntk,dynamicntk} promotes both non-uniform interpolation and extrapolation along the various dimensions of RoPE. YaRN~\cite{yarn} assigns RoPE dimensions into three categories based on frequency, applying extrapolation, NTK, and linear interpolation respectively within each group. Nonetheless, positional embedding in the Transformer model embodies \emph{complex non-uniform information entropy}. Existing methodologies tend to overlook this intricate non-uniformity, which results in the loss of information and constrains the maximum achievable context window size.

Section~\ref{sec:analysis} uncovers two principal empirical insights: \textbf{(1)} Robust positional interpolation necessitates attention to two specific forms of non-uniformity: both in terms of RoPE dimensions and the distribution of token positions. Reduced interpolation suffices for lower RoPE dimensions as well as for earlier token positions, though the precise optimal settings are contingent on the intended extended context length.  
\textbf{(2)} When these aspects of non-uniformity are properly integrated into positional interpolation, the integrity of information within the original RoPE is largely preserved, especially across critical dimensions and early token locations. This practice helps mitigate information loss arising from interpolation, enabling superior initial weights for subsequent fine-tuning. Furthermore, such methods make it possible to achieve up to 8$\times$ context length extension even without fine-tuning.

Inspired by these results, we propose \textbf{LongRoPE}, a robust approach capable of scaling the context window in LLMs to exceed 2 \emph{million} tokens. The design of LongRoPE is anchored in three principal innovations. Firstly, {LongRoPE} capitalizes on multidimensional non-uniform positional interpolation. It determines optimal rescale parameters for RoPE's rotational angles across each dimension of RoPE, contingent upon the tokens’ locations. Given that the process of finding appropriate rescale factors becomes exponentially more complex as the desired extension ratio increases, {LongRoPE} leverages an evolutionary search algorithm enhanced with two distinct optimization strategies to improve search performance. An illustration of the optimized rescaled RoPE can be seen in Fig.~\ref{fig:overview}.

Subsequently, {LongRoPE} implements an efficient, stepwise extension method that enables the model to handle a 2048k context window without requiring direct fine-tuning on ultra-long texts, which are rare and often inaccessible. The approach initiates by identifying a 256k token context length that can be managed by the pre-trained LLM, followed by fine-tuning the model at this intermediate length. Given that the non-uniform positional interpolation facilitates up to an 8$\times$ increase in context size without further fine-tuning, a subsequent search is performed to determine new RoPE rescaling factors on this fine-tuned, length-extended model. This process ultimately expands the context window to 2048k tokens for models such as LLaMA2 and Mistral~\cite{mistral}.

In order to prevent a decline in performance when dealing with the original, shorter context windows, {LongRoPE} maintains ongoing adjustments to the RoPE rescale factors in the expanded LLM. Following the methodology used to increase context length from 256k to 2048k, we also perform down-scaling to 4k and 8k context windows on the LLM that has been fine-tuned at 256k. This is achieved through our search algorithm, which is designed to reduce positional interpolation. At inference time, if the input sequence contains fewer than 8k tokens, we revise RoPE by applying the rescale factors determined from our search process.

Comprehensive testing involving multiple LLMs and a range of long-context challenges provides strong evidence for the efficacy of our approach. Our findings reveal that LongRoPE consistently sustains low perplexity values when evaluated across context lengths spanning from 4k up to 2048k. In addition, it achieves more than 90\% accuracy in passkey retrieval tasks and performs at a similar level to baselines on established benchmarks within the 4096-token context range. The LongRoPE framework is compatible with any LLM employing RoPE embeddings. We intend to make both our code and LongRoPE-2048k models publicly available.

\section{Non-uniformity in Positional Interpolation}
\label{sec:analysis}

\subsection{Preliminary}

Transformer models require explicit positional information, often in the form of position embedding, to represent the order of input tokens. Our work focuses on the RoPE~\cite{rope} position embedding, which is widely used in recent  LLMs. For a token at position index $n$, its corresponding RoPE encoding can be simplified as follows:

\begin{equation}
		\fontsize{7.8}{7.8} \selectfont
	\label{eq:rope}
	\left[ cos(n\theta_0), sin(n\theta_0), cos(n\theta_1), \cdot\cdot\cdot, cos(n\theta_{d/2-1}), sin(n\theta_{d/2-1}) \right]   
\end{equation}

Here, $d$ denotes the size of the embedding space, 
$n\theta_i$ corresponds to the rotary angle assigned to the token located at position $n$, 
and $\theta_i=\theta^{-2i/d}$ specifies the rotation frequency associated with each embedding index. The typical choice for the base value $\theta$ in RoPE is set to 10000.

\textbf{\emph{Context window extension ratio $s$} and positional interpolation}. We introduce $s$ as the proportion between the length of the expanded context window $L'$ and the original window $L$, given by $s=\frac{L'}{L}$.

To extend the context window from $L$ to $L'$, current positional interpolation methods suggest downscaling rotation frequencies $\theta_i$ based on extension ratio $s$. Let $\beta=\theta^{2/d}$, and $\lambda$ denote the actual rescale factor related to $s$, we   unify these positional interpolation methods as follows:

\begin{equation}	
	\fontsize{7.5}{7.5} \selectfont
	\label{eq:unified_rope}
	\begin{aligned}
		\left[cos\left(\frac{n}{\lambda(\beta)^0}\right), sin \left(\frac{n}{\lambda(\beta)^0}\right), cos\left(\frac{n}{\lambda(\beta)^1}\right), \cdot\cdot\cdot,  sin \left(\frac{n}{\lambda(\beta)^{d/2-1}}\right) \right]   
	\end{aligned}
\end{equation}

\textbf{Linear positional interpolation (PI)}. PI~\cite{pi} employs a linear scaling of position indices, constrained by the maximum pre-trained sequence length. When extending to a new target ratio $s$, the angular rotations corresponding to each position are uniformly scaled down by a factor of $\lambda=s$ for every RoPE dimension. While straightforward, this approach compresses the positional representation, causing the positions to be densely packed and thus reducing the model’s capacity to differentiate between tokens that are nearby. Consequently, at larger extension ratios, PI's effectiveness is typically limited.

\textbf{NTK-based interpolation and extrapolation}. Studies such as~\cite{ntk,dynamicntk} examine RoPE through the framework of information encoding, leveraging the Neural Tangent Kernel (NTK) theory~\cite{jacot2018neural,tancik2020fourier}. To address the issue of position congestion found in PI, they propose redistributing the burden of interpolation across the various RoPE dimensions. In this scheme, dimensions associated with lower frequencies receive smaller scaling, while higher-frequency dimensions are scaled more, thus achieving both interpolation and extrapolation over position, described as $\lambda=s^{i}$. An enhanced dynamic NTK variant~\cite{dynamicntk} further modifies the extension ratio at every position, conditional on the sequence's current length. Whereas PI approaches generally require fine-tuning for window extension, NTK-based approaches can enlarge the context window without the need for fine-tuning, albeit typically capped at a 4$\times$ maximum extension ratio.

\textbf{YaRN}~\cite{yarn} offers a notable enhancement in the effectiveness of positional interpolation. The method segments RoPE dimensions into three groups based on their associated frequencies, applying distinct interpolation methods to each. For the highest frequency dimensions, extrapolation is used ($\lambda$=1), while the lowest frequency dimensions are interpolated linearly (PI). The dimensions of intermediate frequency make use of NTK interpolation. A central aspect of YaRN is this categorization of RoPE dimensions; however, the grouping is currently determined through human empirical observation, which can potentially lead to non-optimal results when adapting to new LLMs.

\subsection{Study on Non-uniform Positional Interpolation}

Drawing on insights from NTK and YaRN, we observe that their improvements stem from the use of non-linear mechanisms, particularly by treating frequencies along RoPE dimensions separately to facilitate tailored interpolation and extrapolation. Yet, the non-linear functions presently employed are predominantly based on handcrafted heuristics. This observation prompts two core inquiries: (1) Does the existing approach to positional interpolation represent the best possible method? (2) Could alternative, yet-unexplored non-linearities exist?

To investigate these issues, we employ evolution search (refer to Sec.~\ref{sec:method}) to identify improved non-uniform positional interpolations tailored for LLaMA2-7B. Our approach leverages perplexity as the optimization objective and evaluates each candidate using 5 randomly selected excerpts from the PG19~\cite{pg19} validation set. This empirical study yields several important insights, as detailed below.

\textbf{Finding 1}: \textit{RoPE dimensions exhibit substantial non-uniformities, which are not effectively handled by current positional interpolation methods.}

We determine the optimal value of $\lambda$ for each RoPE dimension as described in Eq.~\ref{eq:unified_rope}. Table~\ref{tbl:exp1} reports the perplexity scores of LLaMA2-7B on the PG19 and Proof-pile~\cite{proof-pile} test datasets using various approaches, all evaluated without fine-tuning. Our method yields pronounced improvements, indicating that both the standard linear interpolation (PI) and the non-uniform techniques (Dynamic-NTK and YaRN) are not optimal. Interestingly, YaRN achieves lower performance compared to PI and NTK on PG19, likely because it does not fully cover the target context window when models are not fine-tuned. Specifically, YaRN's perplexity rises sharply after 7k tokens in an 8k context setting.

Our investigation reveals that the rescaling parameters $\lambda$ in Eq.~\ref{eq:unified_rope} are heterogeneous, contrasting with the constant scaling factor $s$ utilized in PI, the NTK approach, and the group-based method in YaRN. This variability in rescaling considerably boosts the performance of LLaMA2 on language modeling tasks—measured by perplexity—for extended context windows of 8k and 16k tokens, all achieved without additional fine-tuning. The improvement stems from the positional embedding, which retains the original RoPE structure, particularly within key components, thereby facilitating the model's ability to discriminate between nearby token positions.

\textbf{Finding 2}: \textit{RoPE for the initial tokens in the input sequence should be extrapolated with less interpolation.} 

We propose that RoPE should minimize interpolation for the first $\hat{n}$ tokens in input sequences. This stems from the fact that these tokens typically receive substantial attention scores and thus play a pivotal role in attention layers, as reported in Streaming LLM~\cite{streamingllm} and LM-Infinite~\cite{han2023lminfinite}. To test this hypothesis, we expand the context window to 8k and 16k using PI and NTK methods, deliberately omitting interpolation for the initial $\hat n$ tokens (0, 2, ..., 256). If $\hat n = 0$, the procedure defaults to the standard PI and NTK configuration. Table~\ref{tbl:tokenposition} reveals two main findings: \textbf{(1)} exempting the initial tokens from position interpolation notably enhances performance, and \textbf{(2)} the ideal number of non-interpolated starting tokens, $\hat n$, varies depending on the desired context window extension.

\textbf{Finding 3}: \textit{Non-uniform positional interpolation effectively extends LLM context window in both fine-tuning and non-fine-tuning settings.} 

Although our experiments demonstrate that the non-uniform position interpolation we identified substantially enhances extension performance at 8k and 16k context lengths in a zero-shot setting, pushing to even larger context lengths necessitates model fine-tuning. To this end, we perform fine-tuning on LLaMA2-7B using the discovered RoPE configuration tailored for a 64k context window (refer to Appendix for hyperparameter details). Results reported in Table~\ref{tbl:finetune} reveal that our approach consistently outperforms both PI and YaRN, for both pre-trained and fine-tuned LLaMA2-7B. We attribute these gains to leveraging non-uniform positional interpolation, which helps to mitigate information degradation and delivers superior initialization for subsequent fine-tuning.

\textbf{Summary}. This work identifies two key sources of non-uniformity: the RoPE dimensions and token positions. Properly exploiting these in positional interpolation yields marked gains in LLM context length extension performance.

\section{LongRoPE}

\label{sec:method}
In response to these discoveries, we propose LongRoPE, which initially integrates a robust search method designed to leverage both types of non-uniformities. Subsequently, this technique is employed to expand the context window of LLMs past 2 million tokens.

\subsection{Problem Formulation}

These two sources of non-uniformity greatly expand the range of possible solutions and add layers of complexity to the optimization process. To manage this, we conceptualize the multidimensional non-uniform position interpolation optimization as a search problem.

For a LLM targeting a  context window size of $L'$ and lengthy input documents $\mathbf{X}$, where each $\mathbf{x}\in\mathbf{X}$ surpasses $L'$ in token length, we denote the original rotary angle of the $i^{th}$ dimension in RoPE embedding at token position $n$ as  $\frac{n}{\beta^i}$. The optimization problem is then formulated as follows:

\begin{equation}
\begin{gathered}
		\label{eq:problem}

\underset {\mathbf{x}\in\mathbf{X}; \, |\mathbf{x}|\ge L'}{ \operatorname {arg\,min} } \,  \mathcal{L}\, \left( \text{LLM} (\text{RoPE}, \mathbf{X})\right),	\text{where}
\\
\scalebox{0.93}{
$\underset {\begin{subarray}{l} i=0,\cdot\cdot,\frac{d}{2}-1;\\ n\in[ 0,|\mathbf{x}|);\end{subarray}}{\text{RoPE}_(n)}= {\left[\cdot\cdot,cos\left(\mathbb{I}(\hat{\lambda}_{i}, \hat{n})\times\frac{n}{\beta^i}\right),  sin\left(\mathbb{I}(\hat{\lambda}_{i}, \hat{n})\times\frac{n}{\beta^i}\right),\cdot\cdot\right] }$}\\
	\text{where} \;\mathbb{I}(\hat{\lambda}_{i},\hat{n})=\begin{cases}
	1&\mbox{\;  $n<\hat{n}$}\\
{\frac{1}{\lambda}_{i}}&\mbox{\; $n\geq\hat{n}$}
\end{cases}
	\end{gathered}
\end{equation}

In this formulation, we identify the optimal $\mathbf{x}$ from the set $\mathbf{X}$ subject to the constraint $|\mathbf{x}| \ge L'$ that minimizes the objective function $\mathcal{L}$, which depends on the output of the LLM model employing the RoPE positional encoding mechanism. The RoPE at position $n$ is constructed, for indices $i$ ranging from $0$ to $\frac{d}{2}-1$ and token positions $n$ in $[0, |\mathbf{x}|)$, as follows:
\\
$\text{RoPE}_n$ consists of the components $\ldots, \cos\big(\mathbb{I}(\hat{\lambda}_{i}, \hat{n})\cdot \frac{n}{\beta^i}\big), \sin\big(\mathbb{I}(\hat{\lambda}_{i}, \hat{n})\cdot \frac{n}{\beta^i}\big), \ldots$. Here, the indicator function $\mathbb{I}(\hat{\lambda}_{i}, \hat{n})$ is defined by:
\[
\mathbb{I}(\hat{\lambda}_{i},\hat{n}) =
\begin{cases}
1 &\text{if}\; n<\hat{n}\\
{\frac{1}{\lambda}_{i}} &\text{if}\; n\geq\hat{n}
\end{cases}
\]
This parameterization allows for variable scaling based on the token position relative to a cutoff $\hat{n}$, incorporating either unity or a rescaled frequency parameter.

Here, we define a collection of rescaling coefficients, $\mathbb{I}(\hat{\lambda}_{i},\hat{n})$, in order to address both types of non-uniformities present. The variables $\hat{\lambda}_i$ and $\hat{n}$ correspond to the irregularities within the RoPE dimensions and token position indices, respectively. In particular, $\mathbb{I}(\hat{\lambda}_{i},\hat{n})$ serves to adjust the rotational angle associated with the $i^{th}$ dimension in RoPE, where $\hat\lambda_{i}$ functions as the scaling coefficient, and $\hat{n}$ establishes a threshold for token positions. During the first $\hat{n}-1$ token slots, the scaling coefficient $\hat\lambda_{i}$ is not applied; instead, the original RoPE angle $\frac{n}{\beta^i}$ is retained. Once tokens reach positions $n \geq \hat{n}$, the scaling factor $\hat\lambda_{i}$ is utilized.

With a desired context window length of $L'$, the task involves determining the best set of rescale factors ($\mathbb{I}(\hat{\lambda}_{0},\hat{n})$, $\mathbb{I}(\hat{\lambda}_{1},\hat{n})$, ..., $\mathbb{I}(\hat{\lambda}_{i},\hat{n})$, ...) corresponding to each of the $1^{st}$ through $d^{th}$ RoPE dimensions. By applying these rescaled $\text{RoPE}$ parameters to the target $\text{LLM}$, the model aims to minimize the next token prediction loss, denoted as $\mathcal{L}$ (which reflects the perplexity), across input sequences $\mathbf{X}$ that have token lengths of $L'$.

\subsection{Searching the Non-uniform Position Interpolation}

To address the challenge presented in Eq.~\ref{eq:problem}, we propose a straightforward but powerful approach that identifies the best RoPE rescale parameters to maximize the advantages offered by multidimensional non-uniform characteristics in positional embedding.

\textbf{Search space}. We construct an extensive search space that accommodates both non-uniformities. Table~\ref{tbl:searchspace} provides an overview of this search space. In our approach, we permit searching for a distinct rescale factor for each dimension within the RoPE embedding. To streamline the search space, we opt to search over $\lambda_i$ and $\hat{n}$, rather than directly searching $\mathbb{I}(\hat{\lambda}_{i},\hat{n})$, where $\hat{\lambda}_{i}$ is defined as $1/\lambda_i$. 
According to Table~\ref{tbl:searchspace}, the range for $\lambda_i$ extends from a minimum of 1.0 (representing direct extrapolation) up to a maximum of $s\times1.25$ (reflecting greater interpolation than PI), incremented by 0.01, with $s$ denoting the desired ratio of context window expansion.

The parameter $\hat{n}$ specifies how many of the initial token positions are preserved with their native RoPE embedding, foregoing any positional interpolation. In practice, we search for the optimal value of $\hat{n}$ within the set \{0, 1, 2, 4, 8, 12, 16, 20, 24, 28, 32, 64, 128, 256\}. Notably, when $\hat{n}=0$, the rescale factors identified through the search are applied uniformly across all token positions.

\textbf{Evolution-based search}. The search space outlined in Table~\ref{tbl:searchspace} encompasses a vast array of potential positional interpolation solutions, creating substantial barriers to effective exploration. For instance, when applying a $s=4\times$ extension, the number of possible configurations reaches $400^{128/2}\times14$=4$\times10^{167}$, reflecting the exponential growth of the search space as the extension ratio increases. To tackle this complexity, we adopt an evolution search strategy~\cite{guo2020single} and implement two optimization methods to significantly enhance the efficiency of the search process. The overall workflow is summarized in Algorithm~\ref{alg:search}.

\textit{Enhanced initialization of the population}. Rather than solely relying on random generation for the initial set of $P$ rescale factors, we explicitly include the three RoPE rescale factors associated with PI, NTK, and YaRN within the initial group. To complete the population, the remaining $P$-3 members are derived by applying random mutations to these three rescale factors, where each mutation occurs with probability $p$.

\textit{Monotonically non-decreasing constraint}. Following the creation of the initial population, we evaluate the LLM perplexity for every individual. This process involves adjusting the target LLM using the associated RoPE rescale factors, then measuring the perplexity for the input $\mathbf{X}$. The individuals with the lowest perplexities (top-k) are selected as parents for the next phase of evolution. Due to the expansive nature of the search space, simple mutation and crossover methods may inadvertently lead to suboptimal candidates, necessitating additional perplexity calculations. Such redundant evaluations are especially inefficient when $L'$ is large, as each inference for perplexity demands considerable computational resources.

To tackle this issue, we enforce a monotonicity constraint such that the rescaled RoPE factors are non-decreasing, specifically $\lambda_i \leq \lambda_{i+1}$. Only those RoPE configurations that adhere to this constraint are utilized for perplexity assessment on the LLM, which substantially reduces the computational burden of searching. In practice, this constraint mandates that $\lambda_i$ values progressively increase across the RoPE dimensions ($i = 0, \ldots, 63$). The rationale for enforcing monotonicity along these dimensions stems from Neural Tangent Kernel (NTK) theory~\cite{jacot2018neural,tancik2020fourier,ntk}: dimensions corresponding to higher frequencies (lower indices) necessitate less interpolation, resulting in smaller $\lambda_i$, whereas dimensions with lower frequencies (higher indices) can accommodate greater interpolation, thus larger $\lambda_i$.

\begin{algorithm}[t]
	\small
	\caption{The search algorithm}
	\textbf{Input:} target LLM, input samples $\mathbf{X}$, population size $P$, mutation size $N_1$, crossover size $N_2$, maximum iterations $\mathcal{T}$, mutation probability $p$\\

	\begin{algorithmic}[1]
		\label{alg:search}
		\STATE Top-k $\leftarrow$ $\phi$;
		\STATE $\text{P}_0$ $\leftarrow$ \textit{Initialize\_population\_with\_optimization}($P$, $\mathbf{X}$, $p$);
		\FOR{$i = 1$ to $\mathcal{T}$}
		\STATE \textit{Compute\_perplexity} (LLM, $\text{P}_{i-1}$, $\mathbf{X}$);
		\STATE Top-k $\leftarrow$ \textit{Update\_Topk} (Top-k, $\text{P}_{i-1}$);
		\STATE $\text{P}_{mutation}$ $\leftarrow$ \textit{Mutation\_with\_mono\_constraint} (Top-k, $N_1$, $p$);
		\STATE $\text{P}_{crossover}$ $\leftarrow$ \textit{Crossover\_with\_mono\_constraint} (Top-k, $N_2$);
		\STATE $\text{P}_i$ $\leftarrow$ $\text{P}_{mutation}$ $\cup$ $\text{P}_{crossover}$ $\cup$ Top-k;
		\ENDFOR
		\STATE Output the entity in Top-k achieving minimum perplexity;
	\end{algorithmic}
\end{algorithm}

\textbf{8$\times$ expansion without fine-tuning}. Through evolutionary search, we efficiently discover non-uniform RoPE rescale factors that retain crucial dimensions and positional encodings, thereby reducing information degradation from interpolation. As illustrated in Fig.\ref{fig:non-ft}, our approach successfully broadens LLaMA2’s context window from 4k to 32k without necessitating additional fine-tuning steps. In comparison, current techniques like PI, non-uniform NTK, and YaRN result in significant increases in perplexity beyond a 2$\times$ window extension.

\subsection{Extending LLM Context Window to 2048K}
\label{sec:extendto2048k}

\textbf{Progressive extension to 2048k}. In this section, we present our approach for expanding the context window of pre-trained LLMs from the standard 4k length to beyond 2048k. Our use of non-uniform positional interpolation has been shown to enable up to an 8$\times$ increase in context length without the need for fine-tuning. However, when aiming for much larger expansions, such as a 512$\times$ increase, fine-tuning becomes indispensable. A potential strategy involves identifying suitable RoPE rescaling factors for the desired 2048k window and subsequently performing fine-tuning. Nevertheless, this process encounters significant obstacles, chiefly the immense computational resources required for training. Additionally, our empirical observations indicate that effective fine-tuning is particularly difficult when pursuing such large-scale context window extensions (refer to Appendix for further details).

Fortunately, LongRoPE demonstrates robust performance across both the unmodified and fine-tuned, extended LLMs. Thus, we propose a streamlined and incremental approach that attains the desired 2048k context size using only 1k fine-tuning iterations, with a training sequence limited to 256k tokens.

$\Diamond$ \textit{Exploring LongRoPE for expanding pre-trained LLMs to 256k context length}. Using LLaMA2 as a representative, we systematically investigate extending the context window to 128k and 256k tokens. At these targets, the context length is increased by factors of 32$\times$ and 64$\times$, accordingly.

$\Diamond$ \textit{Fine-tuning to 256k}. Subsequently, the pre-trained LLM undergoes fine-tuning to reach a context window of 256k tokens. In detail, LLaMA2 is initially fine-tuned for 400 steps employing RoPE rescaling factors set to 128k. Following this, the RoPE rescaled factors are switched to 256k on the resulting checkpoint, and another 600 steps of fine-tuning are carried out. This sequential approach demonstrates greater efficiency compared to direct fine-tuning to 256k.

$\Diamond$ \textit{Expanding fine-tuned extended LLM to 2048k through LongRoPE search}. In the concluding phase, an additional search is applied to the LLM fine-tuned at 256k sequence length. This step achieves a remarkable increase in the context window, reaching 2048k tokens, accomplished without extra rounds of fine-tuning. The overall extension achieved is $512\times$.

\textbf{Performance within the original context window after extension}.  Expanding the context window to an exceptionally large 2048k tokens results in decreased accuracy in the initial context range. This degradation aligns with observations reported in positional interpolation studies~\cite{pi}, where increasing the position embedding dimensions within the original window confines them to a tighter interval, thus impairing the effectiveness of the language model. Specifically, applying a 512$\times$ extension compresses positions within the standard 4k context window, leading to significant embedding overlap.

To address this issue, we conduct an additional evolutionary search on the extended LLM to fine-tune the RoPE rescale factors specifically for shorter context lengths such as 4k and 8k. For these shorter spans, we restrict the upper limit of the searched $\lambda$ values, given the diminished need for positional interpolation. Inference proceeds with the LLM adaptively modulating the relevant RoPE rescale factors.

\section{LongRoPE}

\label{sec:method}
Drawing on these insights, we propose LongRoPE, which incorporates a novel and efficient search strategy designed to leverage both identified non-uniformities. Utilizing this approach, LongRoPE extends the context window of LLMs to exceed 2 million tokens.

\subsection{Problem Formulation}

These two sources of non-uniformity greatly expand the solution space and increase the difficulty of optimization. To manage these challenges, we recast the multidimensional non-uniform position interpolation optimization task as a search problem.

For a LLM targeting a  context window size of $L'$ and lengthy input documents $\mathbf{X}$, where each $\mathbf{x}\in\mathbf{X}$ surpasses $L'$ in token length, we denote the original rotary angle of the $i^{th}$ dimension in RoPE embedding at token position $n$ as  $\frac{n}{\beta^i}$. The optimization problem is then formulated as follows:

\begin{equation}
\begin{gathered}
		\label{eq:problem}

\underset {\mathbf{x}\in\mathbf{X}; \, |\mathbf{x}|\ge L'}{ \operatorname {arg\,min} } \,  \mathcal{L}\, \left( \text{LLM} (\text{RoPE}, \mathbf{X})\right), \text{where} 
\\
\scalebox{0.93}{
$\underset {\begin{subarray}{l} i=0,\ldots,\frac{d}{2}-1;\\ n\in[ 0,|\mathbf{x}|);\end{subarray}}{\text{RoPE}_(n)}= {\left[\ldots,\cos\left(\mathbb{I}(\hat{\lambda}_{i}, \hat{n})\times\frac{n}{\beta^i}\right),  \sin\left(\mathbb{I}(\hat{\lambda}_{i}, \hat{n})\times\frac{n}{\beta^i}\right),\ldots\right] }$}\\
	\text{where} \quad \mathbb{I}(\hat{\lambda}_{i},\hat{n}) = 
	\begin{cases}
	1 & \text{if} \quad n < \hat{n}\\
	\frac{1}{\lambda}_{i} & \text{if} \quad n \geq \hat{n}
	\end{cases}
	\end{gathered}
\end{equation}

Here, we define a collection of rescaling factors, $\mathbb{I}(\hat{\lambda}_{i},\hat{n})$, which account for both types of non-uniformity in the system. The variables $\hat{\lambda}_i$ and $\hat{n}$ correspond to non-uniformity associated with the RoPE dimensions and token positions, respectively. The function $\mathbb{I}(\hat{\lambda}_{i},\hat{n})$ is used to modify the rotation angle for the $i^{th}$ RoPE dimension by applying the rescale factor $\hat{\lambda}_i$, dependent on the token position threshold $\hat{n}$. For the initial token positions, specifically those less than $\hat{n}$, the rescale factor $\hat{\lambda}_i$ is not utilized; instead, the original RoPE rotary angle $\frac{n}{\beta^i}$ remains unchanged. Once the token positions satisfy $n\ge\hat{n}$, the rescale factor becomes operative.

Given a specified context window length of $L'$, the task is to determine the optimal set of rescale coefficients ($\mathbb{I}(\hat{\lambda}_{0},\hat{n})$, $\mathbb{I}(\hat{\lambda}_{1},\hat{n})$, ..., $\mathbb{I}(\hat{\lambda}_{i},\hat{n})$, ...) across the $1^{st}$ through $d^{th}$ RoPE dimensions. With these rescaling parameters applied to the $\text{RoPE}$ of the target $\text{LLM}$, the model is expected to yield the lowest possible loss for next-token prediction, $\mathcal{L}$ (interpreted as perplexity), when processing input sequences $\mathbf{X}$ of token length $L'$.

\subsection{Searching the Non-uniform Position Interpolation}

To address the challenge posed in Eq.~\ref{eq:problem}, we propose a straightforward and robust technique that identifies ideal RoPE rescale parameters, thereby leveraging the complex multidimensional variations present in position embeddings.

\textbf{Search space}. Our approach incorporates an expansive search space that accommodates both non-uniformities. The configuration of this search space is detailed in Table~\ref{tbl:searchspace}. In particular, we enable the optimization of a unique rescale factor for each individual dimension within the RoPE embedding. To streamline the construction of the search space, we opt to search over $\lambda_i$ and $\hat{n}$ rather than $\mathbb{I}(\hat{\lambda}_{i},\hat{n})$, where $\hat{\lambda}_{i}$ is defined as $1/\lambda_i$. 
Table~\ref{tbl:searchspace} outlines that $\lambda_i$ ranges from a lower bound of 1.0, which corresponds to straightforward extrapolation, up to a maximum value of $s\times1.25$, representing more extensive interpolation relative to PI, with increments of 0.01. Here, $s$ denotes the desired ratio for extending the context window.

The parameter $\hat{n}$ determines how many of the early token positions keep their original RoPE embedding, foregoing position interpolation. In our experiments, we permit $\hat{n}$ to take values from the set \{0, 1, 2, 4, 8, 12, 16, 20, 24, 28, 32, 64, 128, 256\}. If $\hat{n}=0$, then every token position employs the rescale factors identified through the search process.

\textbf{Evolution-based search}. The search space detailed in Table~\ref{tbl:searchspace} encompasses a vast array of positional interpolation configurations, making it exceptionally difficult to traverse efficiently. For instance, with an extension ratio of $s=4\times$, the number of possible combinations amounts to $400^{128/2}\times14 = 4\times10^{167}$. Increasing the extension ratio further causes the search space to grow at an exponential rate. To tackle this challenge, we employ evolution search~\cite{guo2020single} and integrate two optimization strategies that substantially enhance the effectiveness of the search process. The comprehensive workflow of our search is outlined in Algorithm~\ref{alg:search}.

\textit{Optimized initial population generation}. In place of a fully random initialization for the population of $P$ rescale factors, we explicitly incorporate the three RoPE rescale factors related to PI, NTK, and YaRN as members of the starting population. The other $P$-3 individuals are then obtained by applying random mutations to these three rescale factors, each with a mutation probability of $p$.

\textit{Monotonically non-decreasing constraint}. Once the initial population has been produced, we evaluate LLM perplexity for each candidate. This involves applying the designated RoPE rescale factors to the target LLM and measuring the perplexity on input $\mathbf{X}$. The most promising $k$ candidates are selected to serve as parents for the evolutionary process. Nonetheless, due to the broadness of the search space, basic mutation and crossover operations may frequently investigate inferior options, thereby prompting redundant perplexity assessments. This inefficiency becomes more pronounced when $L'$ is large, since each perplexity evaluation requires considerable inference time.

To mitigate this issue, we enforce a constraint of monotonic non-decreasing order on the sampled RoPE rescaling factors, such that $\lambda_i \leq \lambda_{i+1}$. Only those RoPE configurations meeting this requirement are considered during LLM perplexity testing, which greatly limits the computational burden of the search process. In particular, we stipulate that the values of $\lambda_i$ must increase monotonically across RoPE dimensions indexed by $i$ (where $i=0,\ldots,63$). The rationale for this monotonic progression stems from NTK theory~\cite{jacot2018neural,tancik2020fourier,ntk}, which posits that lower dimensions—corresponding to higher frequencies—necessitate less interpolation (hence a smaller $\lambda_i$), while higher dimensions—associated with lower frequencies—support greater interpolation (thus a larger $\lambda_i$).

\begin{algorithm}[t]
	\small
	\caption{The search algorithm}
	\textbf{Input:} target LLM, input samples $\mathbf{X}$,   population size $P$, mutation size $N_1$, crossover size $N_2$, max iterations $\mathcal{T}$, mutate probability $p$\\

	\begin{algorithmic}[1]
		\label{alg:search}
		\STATE Top-k$\leftarrow\phi$;	
		\STATE $\text{P}_0$ $\leftarrow$ \textit{Initialize\_population\_with\_optimization} ($P$, $\mathbf{X}$, $p$); 
		\FOR{$i=1$ to $\mathcal{T}$}
		\STATE \textit{Compute\_perplexity} (LLM, $\text{P}_{i-1}$, $\mathbf{X}$);
		\STATE Top-k $\leftarrow$ \textit{Update\_Topk} (Top-k, $\text{P}_{i-1}$);
		\STATE$\text{P}_{mutation}$ $\leftarrow$ \textit{Mutation\_with\_mono\_constraint} (Top-k, $N_1$, $p$); 
		\STATE $\text{P}_{crossover}$ $\leftarrow$ \textit{Crossover\_with\_mono\_constraint} (Top-k, $N_2$);
		\STATE $\text{P}_i$ = $\text{P}_{mutation}$ $\cup$ $\text{P}_{crossover}$ $\cup$ Top-k;
		\ENDFOR
		\STATE Output the individual exhibiting the minimum perplexity found in Top-k;
	\end{algorithmic}
\end{algorithm}

\textbf{8$\times$ extension achieved without fine-tuning}. Leveraging our evolutionary search approach, we successfully determine non-uniform RoPE rescale factors that strategically protect critical dimensions and positions, thereby reducing information loss resulting from interpolation. As illustrated in Fig.\ref{fig:non-ft}, this enables our method to scale the context window of LLaMA2 from 4k to 32k tokens absent any fine-tuning procedure. On the other hand, alternative approaches including PI, as well as non-uniform NTK and YaRN, result in significantly increased perplexity beyond a 2$\times$ context window expansion.

\subsection{Extending LLM Context Window to 2048K}
\label{sec:extendto2048k}

\textbf{Progressive extension to 2048k}. In this section, we present our approach for expanding the context window of pre-trained LLMs from the standard 4k tokens up to more than 2048k tokens. Our results show that employing non-uniform positional interpolation allows for an 8$\times$ extension without the need for fine-tuning. However, achieving substantially larger extensions (for example, up to 512$\times$) necessitates fine-tuning. A potential strategy involves determining RoPE rescaled factors tailored to the desired 2048k window, followed by fine-tuning. Nevertheless, this approach encounters significant obstacles owing to the extremely high computational costs involved in training. In addition, our observations indicate that it is quite difficult to effectively fine-tune LLMs when targeting such substantial extension ratios (refer to the Appendix for details).

LongRoPE demonstrates efficacy when applied to both the baseline and adapted, extended LLMs. In light of this, we present a streamlined, incremental approach that attains the intended 2048k context window using merely 1k fine-tuning steps, restricting the training sequence length to 256k.

$\Diamond$ \textit{LongRoPE search for scaling pre-trained LLMs to 256k}. Using LLaMA2 as a representative case, we explore configurations aiming for context window lengths of 128k and 256k tokens. In this process, the respective expansion rates are 32$\times$ and 64$\times$.

$\Diamond$ \textit{Fine-tuning to 256k}. Next, we adapt the pre-trained LLM to handle a 256k context window. To do so, LLaMA2 is initially fine-tuned for 400 steps employing RoPE rescaled factors targeted at 128k. After completing this phase, the RoPE rescaled factors are swapped to those for 256k, and a further 600 steps of fine-tuning are performed starting from the resulting checkpoint. This incremental approach demonstrates greater efficiency compared to attempting direct fine-tuning for the 256k context size.

$\Diamond$ \textit{Adapting fine-tuned extended LLM to a 2048k context window via LongRoPE search}. In the final stage, we carry out a secondary search step using the 256k fine-tuned LLM, enabling a substantial expansion to a 2048k context window without requiring any additional fine-tuning. This process yields a cumulative extension factor of $512\times$.

\textbf{Restoring performance in shorter context windows}.  When scaling up to a vast 2048k context window, we observe a degradation in performance on the initial context window. This phenomenon has been previously identified with positional interpolation~\cite{pi}, which compresses the position embeddings of the original window into a significantly smaller range when embedding dimensions are mapped to support longer contexts, resulting in diminished model effectiveness. In particular, applying a 512$\times$ scaling factor leads to substantial overcrowding of positional representations within the standard 4k window.

To address this issue, we conduct an additional evolution search on the expanded LLM, specifically tuning RoPE rescale factors for shorter context lengths such as 4k and 8k. For these shorter spans, we limit the maximum searched $\lambda$ since less positional interpolation is necessary. While running inference, the LLM modifies the relevant RoPE rescale factors in real time.

 \section{Experiments}

\subsection{Setup}
\textbf{Evaluation Tasks and models}. LongRoPE is integrated with LLaMA2-7B and Mistral-7B, and their effectiveness is assessed across three dimensions: (1) perplexity scores for LLMs with extended contexts when processing lengthy documents; (2) a Passkey retrieval challenge, which tests the model’s capability to extract a straightforward passkey amid largely unrelated text; and (3) established LLM benchmark tasks constrained to a 4096-token context window.

\textbf{Fine-tuning}. For the LLaMA2 model, a learning rate of $2 \times 10^{-5}$ is employed with linear scheduling and a global batch size of 32. The model undergoes 400 steps of fine-tuning on the Redpajama~\cite{redpajama} dataset, which is divided into 128k-length segments and delimited by BOS and EOS tokens. Subsequently, using the completed checkpoint, a further 600 steps are carried out in order to expand the context window to 256k. The training for the 128k context size utilizes 8 A100 GPUs along with the distributed training platform~\cite{cube}; scaling to 256k context window requires 16 A100 GPUs. For the Mistral model, fine-tuning proceeds with a fixed learning rate of $1 \times 10^{-6}$ and a global batch size of 64. Both the 128k and 256k context models follow the procedures from YaRN~\cite{yarn}, performing 400 steps on Together Computer’s Long-Data Collections~\cite{mistraldata} with a 16k sequence length and utilizing 4 A100 GPUs throughout training.

\textbf{Search}. For target window sizes up to 256k, our configuration uses $P$=64, $N_1$=$N_2$=16, $p$=0.3, and $\mathcal{T}$=40, choosing the top 32 candidates for mutation and crossover operations in each iteration. Perplexity is assessed based on five randomly selected samples from the PG19 validation dataset, ensuring each meets the minimum required target context length. When the target window exceeds 512k, we reduce the population, mutation, and crossover parameters by half. In this case, perplexity evaluation utilizes three randomly sampled validation documents from Pile-Books3~\cite{pile}.

\textbf{Baselines}. In order to extend the context window to 2048k, we performed fine-tuning on models originally set at 128k and 256k context lengths. Consequently, LongRoPE-2048k (ft=128k) corresponds to LLaMA2, while LongRoPE-2048k (ft=256k) refers to Mistral. These four models are evaluated alongside the leading approaches for expanding context windows, focusing on open-source LLMs refined through positional interpolation using techniques such as PI, NTK, and YaRN. The comparison includes models like Together-32k~\cite{together}, Code LLaMA~\cite{codellama}, LongLoRA-full-FT-100k~\cite{longlora}, YaRN-LLaMA, and YaRN-Mistral~\cite{yarn}.

\subsection{Main Results}

\label{sec:mainresult}
\textbf{Language modeling over long sequences up to 256k tokens}. Our initial experiments focus on benchmarking against leading long-context LLMs at the 256k evaluation length. To assess generalization capabilities, we perform evaluations on two datasets: the Proof-pile~\cite{pg19} and PG19~\cite{pile} test sets. Perplexity is measured at multiple context sizes utilizing a sliding window approach of 256 tokens. For PG19, evaluation is conducted over the entire set of 100 test documents. In the case of Proof-pile, following the protocol established by YaRN~\cite{yarn}, we randomly sample 10 instances, each possessing a minimum length of 128k tokens.

Table~\ref{tbl:proof-pile} and Table~\ref{tbl:pg19} present perplexity comparisons between LLaMA2 and Mistral models, each augmented through various interpolation strategies, evaluated on Proof-pile and PG19 datasets. Two main findings emerge: \textbf{(1)} our extended variants demonstrate a consistent reduction in perplexity across evaluation lengths ranging from 4k up to 256k, indicating effective utilization of longer contextual information. \textbf{(2)} Remarkably, despite a substantial context window increase by a factor of 16—an environment where models typically struggle to maintain performance at shorter lengths—our LongRoPE-2048k configurations deliver superior results relative to current state-of-the-art baselines for context windows up to 256k.

\textbf{Long sequence language modeling beyond 2000k}. To assess model performance on exceptionally lengthy texts, we utilize the Books3~\cite{pile} dataset. For the purpose of efficient evaluation, we sample 20 books at random, each containing more than 2048k tokens, and apply a sliding window approach with a size of 256k.

As detailed in Table~\ref{tbl:books3}, LongRoPE effectively broadens the context window of LLaMA2-7B and Mistral-7B models to 2048k, while attaining perplexity results at shorter context lengths (8k-128k) that are on par with, or better than, existing baselines. A considerable disparity in performance emerges between LLaMA2 and Mistral at 2048k. While Mistral surpasses baselines when evaluated on shorter contexts, its perplexity rises above 7 for lengths exceeding 256k. The behavior of LLaMA2 is consistent with anticipated trends: as context length grows, perplexity gradually declines, with only slight increases at 1024k and 2048k. Notably, for LLaMA2, LongRoPE-2048k demonstrates enhanced results when fine-tuned at 256k as opposed to 128k, attributable to the lower secondary extension ratio (8$\times$ compared to 16$\times$). Conversely, Mistral achieves higher performance when fine-tuned on a window of 128k. This can primarily be attributed to the adopted strategy (in accordance with YaRN) for Mistral's 128k and 256k fine-tuning, where a 16k training length is used, influencing the model's subsequent ability to extend its context window after fine-tuning.

\textbf{Passkey retrieval}. Next, we examine the influence of context window size on generative performance. Adhering to the passkey retrieval benchmark introduced by~\cite{passkey}, we assess models using a synthetic protocol. In this setup, a randomly selected five-digit passkey is embedded within an extended document, and the model must extract this hidden value. The structure of the prompt employed is provided in the appendix. For our experiments, we repeat the passkey retrieval process 10 times, each time positioning the passkey at a randomly chosen location uniformly distributed throughout the context window of the evaluation.

Fig.~\ref{fig:passkey} presents retrieval accuracy in comparison to baseline approaches. Notably, the accuracy of previous models deteriorates rapidly, reaching zero for sequences longer than 128k tokens. In contrast, our LongRoPE-LLaMA2-2048k (ft=256k) exhibits robust performance on the demanding task of extracting a passkey across up to two million tokens, sustaining high retrieval accuracy ($\ge$90\%) consistently from 4k through 2048k tokens. Similarly, LongRoPE-Mistral-2048k (ft=128k) maintains perfect (100\%) accuracy up to 1800k, experiencing a decrease to 60\% at 2048k—an outcome in line with the rise in perplexity at this length reported in Table~\ref{tbl:books3}.

\textbf{Evaluation on standard benchmarks in the native context window}. LongRoPE-2048k models are assessed on the original context window through the Hugging Face Open LLM Leaderboard~\cite{huggingfaceboard}, covering both zero-shot and few-shot scenarios. Specifically, we perform 25-shot evaluation on ARC-Challenge~\cite{allenai:arc}, 10-shot testing with HellaSwag~\cite{zellers2019hellaswag}, 5-shot assessment for MMLU~\cite{mmlu}, and zero-shot evaluation on TruthfulQA~\cite{lin2021truthfulqa}.

As indicated in Table~\ref{tbl:commonsense}, our models deliver similar outcomes on the standard benchmark, which was developed for a limited context window. Notably, they exceed the performance of the original Mistral on TruthfulQA by +0.5\%. Although LongRoPE-LLaMA2-2048k—fine-tuned at 256k tokens—experiences somewhat greater declines in performance, its results for the majority of tasks stay well within acceptable limits.

\subsection{Ablation Results}

\textbf{Evaluating the Effectiveness of Second Positional Interpolation}. As part of our progressive extension approach, we employ our search algorithm to apply a second round of non-uniform positional interpolation to already fine-tuned and extended LLMs. To assess its efficacy, we perform experiments using the LLaMA2-256k model after fine-tuning. This model is further expanded to sequence lengths of 512k, 1024k, and 2048k through PI and YaRN methods. Results presented in Table~\ref{tbl:secondsearch} demonstrate that our non-uniform positional interpolation maintains a stable perplexity across increasing extension lengths. By comparison, both PI and YaRN exhibit rapid growth in perplexity as the extension ratio rises.

\textbf{Recovery performance at reduced context lengths.} To address the observed degradation in performance when working with shorter contexts, we recalibrated the RoPE factors for LongRoPE-2048k using our search methodology. In particular, we imposed stricter constraints on the maximum scale factors considered during the search, promoting reduced interpolation for short 4k and 8k contexts. Table~\ref{tbl:shortperformance} presents the perplexity metrics for LongRoPE-LLaMA2-2048k evaluated on the Proof-pile dataset at both 4k and 8k lengths, alongside mean LLM benchmark accuracy. The findings indicate that these adjustments yield marked improvements in performance under short context scenarios.

\textbf{Analysis on the two forms of non-uniformities}. To assess the influence of each type of non-uniformity on model effectiveness, we conducted an ablation study. Two distinct experiments were performed: (i) we scaled LLaMA2-7B to shorter sequence lengths of 16k and 32k using three strategies—PI, adjusting the RoPE dimension alone, and modifying both non-uniformities; (ii) we extended a 256k-length fine-tuned LLaMA2 up to 2048k, adhering to the same experimental framework. Perplexity was measured for all settings without further fine-tuning. As presented in Table~\ref{tbl:searchdimension}, the application of non-uniformity within the RoPE dimension achieves notably lower perplexity compared to linear interpolation methods such as PI. Furthermore, non-uniform adjustments to token position enhance performance at 16k and 32k lengths but appear less effective at 2048k, perhaps due to the challenges posed by such extreme sequence lengths. Relying solely on the initial tokens without interpolation fails to yield meaningful benefits, and we propose to explore this limitation in future research.

\section{Related Works}

Beyond techniques centered on position interpolation, this section reviews literature on alternative methods.

\textbf{Retrieval-based} techniques incorporate an external memory system that stores extensive past information and utilize retrieval mechanisms to access pertinent documents during inference~\cite{longllama,wang2023augmenting,borgeaud2022improving}. Such methods generally require deliberate architectural changes to the underlying LLMs. By comparison, our approach adopts a more efficient and straightforward strategy, involving only limited alterations to positional embeddings. This allows us to support a broader array of tasks involving extended contexts, including long-form document summarization and few-shot learning, surpassing the capabilities offered solely by retrieval methods.

\textbf{Attention-based context window extensions}. In addition to positional embedding interpolation, certain studies extend the input context of LLMs while preserving their original context window by modifying attention mechanisms~\cite{han2023lminfinite, streamingllm,parallelwindow}. Central to this approach is the use of innovative attention masks to address the attention explosion problem associated with new positional tokens. These techniques are compatible with positional interpolation, and can be jointly applied.

\textbf{Fine-tuning based approaches} are concerned with efficiently adapting pre-trained LLMs through modified positional embeddings to accommodate longer sequences. For example, methods such as Code LLaMA~\cite{codellama}, LLaMA2 Long~\cite{xiong2023effective}, and ScaledRoPE~\cite{liu2023scaling} address this by employing a substantially larger base value for RoPE followed by fine-tuning at the intended target sequence length. In contrast, our approach enables adaptation to a broad range of target lengths and is capable of scaling beyond 2M tokens. Additionally, as fine-tuning for extensive context windows (e.g., greater than 128k) is highly resource-intensive on GPUs, approaches like LongLoRA~\cite{longlora} and PoSE~\cite{zhu2023pose} have been introduced to reduce computational demands. Our technique operates independently and can be combined with such efficient fine-tuning methods.

\section{Conclusion}

This paper introduces LongRoPE, a technique that dramatically pushes the context length of LLMs up to 2048k, all while preserving their performance in the original, shorter context range. By leveraging two types of non-uniformity within RoPE positional embeddings via a computationally efficient evolutionary search, LongRoPE delivers dual advantages: it generates effective initializations for subsequent fine-tuning and allows for up to an 8$\times$ expansion of the context window without any fine-tuning. On this foundation, we implement a progressive extension framework that utilizes LLMs fine-tuned at 256k context to achieve a 2048k context length, requiring no additional fine-tuning. Thorough empirical evaluation demonstrates LongRoPE's efficacy. We anticipate that LongRoPE-2048k will catalyze novel applications that depend on long contexts and drive ongoing investigation in this domain.

\section*{Broader Impacts} The objective of this research is to contribute to the progress of Machine Learning as a discipline. While our study may have various implications for society, we do not identify any that necessitate particular mention in this context.

	\balance

\end{document}